% !TeX root = ../main.tex
\chapter*{KẾT LUẬN VÀ HƯỚNG PHÁT TRIỂN}
\addcontentsline{toc}{chapter}{KẾT LUẬN VÀ HƯỚNG PHÁT TRIỂN}
\markboth{KẾT LUẬN VÀ HƯỚNG PHÁT TRIỂN}{KẾT LUẬN VÀ HƯỚNG PHÁT TRIỂN}

\section*{1. Kết quả đạt được}

Trong khuôn khổ chuyên đề, các kết quả chính đã đạt được bao gồm:

\begin{itemize}
  \item Hệ thống hóa cơ sở lý thuyết về xác thực tập trung, giao thức OAuth 2.0
  và OpenID Connect, cũng như các phương pháp triển khai tự động.
  \item Thiết kế kiến trúc ba tầng cho hệ thống định danh tập trung dựa trên
  Keycloak, với điểm nhấn là quản lý cấu hình định danh như mã nguồn.
  \item Triển khai thành công dịch vụ Keycloak cùng PostgreSQL và ứng dụng
  demo, minh họa luồng đăng nhập SSO qua OIDC với PKCE.
  \item Xây dựng được quy trình CI/CD và định hướng triển khai theo mô hình
  GitOps trên Kubernetes (k3s).
  \item Đánh giá hệ thống về chức năng và bảo mật, đề xuất khung đo lường hiệu
  năng cho giai đoạn tiếp theo.
\end{itemize}

Bảng~\ref{tab:muctieu} đối chiếu các mục tiêu đặt ra với kết quả đạt được.

\begin{table}[htbp]
\centering
\caption{Đối chiếu mục tiêu và kết quả đạt được}
\label{tab:muctieu}
\small
\begin{tabular}{p{6cm} p{4.2cm} p{2.6cm}}
\toprule
\textbf{Mục tiêu} & \textbf{Kết quả} & \textbf{Mức độ}\\
\midrule
Xây dựng hệ thống định danh tập trung & Đã triển khai Keycloak, SSO OIDC & Đạt\\
Quản lý cấu hình như mã nguồn & Realm JSON và config-cli & Đạt\\
Tự động hóa triển khai & CI/CD, thiết kế GitOps & Đạt một phần\\
Đánh giá bảo mật & Phân tích STRIDE và biện pháp & Đạt\\
\bottomrule
\end{tabular}
\end{table}

\section*{2. Hạn chế}

Bên cạnh kết quả đạt được, chuyên đề còn một số hạn chế:

\begin{itemize}
  \item Hệ thống mới được kiểm thử ở quy mô nhỏ, chưa đánh giá đầy đủ dưới tải
  cao và trong điều kiện nhiều bản sao.
  \item Một số hạng mục tự động hóa (IaC, GitOps, giám sát) mới dừng ở mức
  thiết kế, chưa hoàn thiện và vận hành ổn định.
  \item Chưa triển khai các cơ chế nâng cao như liên kết định danh (SAML/LDAP),
  zero-trust giữa các dịch vụ, hay đa vùng dự phòng.
  \item Phân tích hiệu năng và bảo mật còn ở mức định tính, cần bổ sung số liệu
  định lượng.
\end{itemize}

\section*{3. Hướng phát triển}

Trên cơ sở các kết quả và hạn chế nêu trên, hướng phát triển thành đồ án tốt
nghiệp bao gồm:

\begin{itemize}
  \item \textbf{Hoàn thiện tự động hóa:} xây dựng đầy đủ IaC (Terraform,
  Ansible), GitOps hoàn chỉnh với ArgoCD và quản lý bí mật bằng Vault/Sealed
  Secrets.
  \item \textbf{Tăng cường tính sẵn sàng cao:} cấu hình nhiều bản sao Keycloak,
  sao chép PostgreSQL, kiểm thử khả năng tự phục hồi và cân bằng tải.
  \item \textbf{Đo lường định lượng:} sử dụng k6/Locust để đo độ trễ, thông
  lượng và MTTR; trực quan hóa bằng Grafana.
  \item \textbf{Mở rộng bảo mật:} triển khai zero-trust với mTLS giữa các dịch
  vụ, liên kết định danh với LDAP/Active Directory và đăng nhập WebAuthn.
  \item \textbf{Mở rộng quy mô:} nghiên cứu triển khai đa vùng và kịch bản dự
  phòng thảm họa (DR) với các chỉ tiêu RTO/RPO.
\end{itemize}

\vspace{0.5cm}
Tóm lại, chuyên đề đã xây dựng được nền tảng vững chắc cho một hệ thống xác
thực tập trung có khả năng triển khai tự động, đồng thời mở ra nhiều hướng
nghiên cứu và phát triển phù hợp với định hướng kỹ thuật hệ thống và bảo mật.
